\section{LEMAITRE: Hubble diagram project}

The LEMAITRE project aims at constructing the largest Hubble diagram by using an independent set of supernovae. One its objectives is to tackle the $(w_0-w_a)$ tension observed 
by recent results from \cite{Brout_2022}, \cite{Rubin_2023} and \cite{DES_2024} supernova analyses as well as the combination with BAO results from \cite{DESI_DR2_2025}. It is 
important to note that the low redshift supernova sample used in all three analyses is the same and that it seems to be the main motivating factor for the tension. 
By analyzing a completely new dataset for LEMAITRE, an independent assessment of this tension will be provided. Additionally with the LEMAITRE project comes a new 
analysis pipeline from end-to-end, starting from the lightcurve extraction and calibration up to the cosmological analysis. My work during the PhD involved the continued 
development and testing of a portion of the LEMAITRE pipeline named NaCl (see Chapter 3 and 4) which aims at training a spectrophotometric model. \\

\subsection{Datasets: ZTF, SNLS, HSC}

There are three datasets used in the LEMAITRE project: \\

\noindent \textbf{\underline{ZTF}} \\

The first will be the ZTF DR2 dataset making up the low-$z$ sample for the LEMAITRE analysis of about 2600 supernova. More details on the ZTF dataset and telescope are explained 
in \ref{sec:ztf}. \\

\noindent \textbf{\underline{SNLS}} \\

The second dataset is the SuperNova Legacy Survey (SNLS)  (Astier et al., 2006), a type Ia supernova survey conducted with the MegaCam (Boulade et al., 2003) observing system. 
The camera is mounted on the Canada-France-Hawaii 3.6m telescope (CFHT) located on top of the Maunakea volcano, Hawaii.
The focal plane of the MegaCam camera is made up of 40 CCDs (model e2v CCD42-90) each having 
2048 $\times$ 4612 pixels. The CFHT is constructed with two mirrors, the primary mirror has a parabolic shape and is made from a glass-ceramic
material which makes it immune to thermal distortion. The camera is placed in the focal plane of the primary mirror. In addition, a Wide Field
Corrector (WFC) is placed in front of the camera to guarantee a large field of view. \\ 

In order to classify observed supernovae and acquire their redshifts, 
a spectroscopic follow-up is implemented. These additional
observations mostly rely on the European Southern Observatory Very Large Telescope (ESOVLT), 
the Gemini-North and South and the Keck-I and II instruments (Balland et al., 2009;
Astier et al., 2006). As the photometric survey observes more type Ia candidates than what
can be spectroscopically observed afterwards, a target selection process must be implemented. 
In practice, a first light-curves fit is performed in real time which allows to eliminate most of the type II contamination in the survey
For the remaining targets, a limit of mi = 24.5mag (i-band magnitude) has been established, over which a supernova is less
likely to be spectroscopically observed.
The MegaCam camera footprint is a 0.96 × 0.94 = 0.90deg2
rectangle. As a consequence, the chosen observation strategy for SNLS is no longer a wide scan of a large fraction of the sky as for ZTF 
but rather consists is observing deeply (up to z  1) in four
dedicated fields equally distributed in right ascension (see figure 3.10 and table 3.1)
The luminosity of transients is collected through 5 photometric filters, g, r, i, i21
and z,
spanning the visible and near-infrared wavelength ranges. Figure 3.11 presents the transmission of the SNLS filters as a function of wavelength. 
As for ZTF earlier, we highlight the
fact that the transmissions are sensor-dependent by displaying the average transmission
over the focal plane and a transmission curve for four different sensors.
In this thesis, the analysed SNLS data consists in the five years survey (SNLS 5y), which
started in the second half of 2003 and ended in 2009. The dataset comprises SNLS 3y data
as well as data from the 4th and 5th years which will be published in Betoule et al. in prep.
Figure 3.12 presents the number of exposures per day as a function of time for the entire
SNLS survey. \\


\noindent \textbf{\underline{HSC}} \\


The Hyper Suprime-Cam Subaru Strategic Program (Aihara et al., 2022), also denoted HSCSSP, is an astronomical 
observation program which allocated some telescope time to the 
observation of type Ia supernovae at the end of the 2010’s. Its primary system is the Hyper
Suprime Cam, a camera mounted on the Subaru 8.2m telescope. As for the CFHT, the
Subaru telescope is located on top of the Maunakea volcano, Hawaii (see figure 3.13).
The focal plane of the Hyper Suprime-Cam camera includes 116 Hamamatsu CCDs 
each totalizing 2048 × 4096 = 8.4Mpx (Hyper Suprime Cam documentation). The Subaru
telescope also follows a Prime Focus architecture. Figure 3.14 presents a side view of the
Hyper Suprime Cam mounted on the Subaru telescope.
The redshifts of the HSC supernovae are determined by measuring the redshift of their
host galaxies. These redshifts were mainly collected by the AAT/AAOmega and DESI
instruments, completed by observations from the Subaru/FOCAS and VLT instruments
(Suzuki et al. in prep). As of today, around 500 redshifts have been measured and observations are still going on, in particular in the XMM field (see figure 3.15).
The HSC-SSP collaboration also developed four methods to determine the redshifts directly from photometry. The Direct Empirical Photometric method (DEmP) and the TPZ
method both rely on machine learning to infer redshifts. The first approach consists in
fitting a polynomial function modeling the redshift as a function of magnitude (Hsieh and
Yee, 2014) while the second one uses random forest algorithms (Carrasco Kind and Brunner,
2013). For complementary measurements, the LePhare and MIZUKI methods both relies on
template fitting with priors. While the LePhare approach uses redshift distributions priors
from the 30-band photometric redshifts of the COSMOS survey (Ilbert et al., 2009; Coupon
et al., 2015), the MIZUKI approach uses Bayesian priors (Tanaka, 2015). 
The Hyper Suprime-Cam camera footprint is a 1.8deg2 disc which is about twice the
size of the MegaCam footprint. Although the main focus of the HSC-SSP is to build a galaxy
catalog for weak lensing analysis, the observation cadence of its deep fields were adapted
for supernovae observation. Because the main goal of the HSC supernova survey is to
observe faint transients at high redshift (up to z  1.6), the chosen observation strategy
is also similar to SNLS. In fact, the supernovae observation is restricted to four dedicated
deep fields (XMM-LSS, E-COSMOS, ELAIS-N1, DEEP2-F3) and two ultra-deep fields (SXDS
and COSMOS), partially overlapping the SNLS D1 and D2 fields respectively. Figure 3.15
presents a summary of the HSC deep and ultra-deep fields locations. 
The luminosity of transients is collected through 7 photometric filters, g, r, r22
, i, i2, z
and y, spanning the visible and near-infrared wavelength ranges. Figure 3.16 presents the
transmission of the HSC filters as a function of wavelength. Although the transmissions are
sensor-dependent as for the ZTF and SNLS filters, we only display the average transmission
curves over the focal plane for clarity.

In this thesis, the analysed HSC data consists in the HSC supernova survey, which
started with a first observation season in 2017 and ended with a second observation season
in 2020. Figure 3.17 presents the number of exposures per day as a function of time for the
entire HSC survey

\subsection{Pipeline overview}
\subsection{Cosmological analyses with LEMAITRE}
\subsection{Uncertainties in SN measurements?}
